% =====================================================================
%  CHƯƠNG 3 — CÁC THẾ HỆ KIẾN TRÚC RAG
% =====================================================================
\chapter[Các thế hệ kiến trúc Retrieval-Augmented Generation]{Các thế hệ kiến trúc\\
Retrieval-Augmented Generation}\label{chap:rag}

RAG không phải một kiến trúc mà là một họ kiến trúc tiến hóa qua ba thế hệ~\cite{gao2023survey}
(Hình~\ref{fig:thehe}). \textbf{Naive RAG} chia văn bản, nhúng, tìm top-$k$ rồi sinh. \textbf{Advanced RAG} thêm
các bước tối ưu trước truy xuất (viết lại câu hỏi) và sau truy xuất (xếp lại, nén ngữ cảnh). \textbf{GraphRAG}
thay đơn vị tri thức từ đoạn văn rời rạc sang đồ thị thực thể và cộng đồng, để trả lời được câu hỏi đa bước và
câu hỏi tổng hợp. Chương này đi qua cả ba thế hệ, và với mỗi thế hệ đo trên chính kho tri thức của bài xem nó cải
thiện được gì, tốn thêm bao nhiêu.

\hinhtikz{t3_generations}{Ba thế hệ RAG xếp tầng: Naive RAG, Advanced RAG với tối ưu tiền và hậu truy xuất,
GraphRAG với đồ thị thực thể và tóm tắt cộng đồng}{thehe}

\section{Naive RAG và các điểm nghẽn}
\subsection{Quy trình hoạt động của Naive RAG}
Naive RAG có ba pha (Hình~\ref{fig:naive}). Pha \textbf{lập chỉ mục} chạy ngoại tuyến: tài liệu được chia thành
đoạn (chunk), mỗi đoạn được mã hóa thành vector $\mathbf{d}_j = f(\text{đoạn}_j) \in \mathbb{R}^{768}$ và lưu vào
chỉ mục. Pha \textbf{truy xuất} chạy khi có câu hỏi: câu hỏi được mã hóa thành $\mathbf{q} = f(\text{câu hỏi})$ và
$k$ đoạn có độ tương đồng cosin lớn nhất được chọn:
\begin{equation}
  \operatorname{sim}(\mathbf{q}, \mathbf{d}_j) = \frac{\mathbf{q}^\top \mathbf{d}_j}{\lVert\mathbf{q}\rVert\,\lVert\mathbf{d}_j\rVert},
  \qquad \mathcal{Z}_k(\mathbf{q}) = \operatorname*{arg\,top-}k_{j}\ \operatorname{sim}(\mathbf{q}, \mathbf{d}_j).
\end{equation}
Pha \textbf{tạo sinh} nối $k$ đoạn vào prompt theo khuôn ``BẰNG CHỨNG: [1]\dots[k] — CÂU HỎI — TRẢ LỜI'' và để LLM
sinh câu trả lời.

\hinhtikz{t3_naive}{Naive RAG ba pha: lập chỉ mục ngoại tuyến, truy xuất và tạo sinh trực tuyến, kèm ba điểm
nghẽn chính}{naive}

\subsection{Các điểm nghẽn của Naive RAG}
Mỗi pha có một điểm nghẽn, và lỗi ở pha trước không pha sau nào sửa được:

\begin{enumerate}
  \item \textbf{Chia đoạn cắt rời ngữ cảnh.} Cửa sổ 400 ký tự có thể cắt giữa câu ``chống chỉ định khi eGFR
        $< 30$''. Đoạn chứa nửa đầu không còn con số, đoạn chứa nửa sau không còn tên thuốc. Mục~\ref{sec:chunking}
        đo hiện tượng này.
  \item \textbf{Truy xuất trượt từ khóa hoặc trượt ngữ nghĩa.} Truy xuất dày hiểu ``tiểu đường'' gần nghĩa ``đái
        tháo đường'' nhưng dễ bỏ qua một mã cụ thể như ``Nghị định 232''; truy xuất thưa thì ngược lại. Với cách
        chia đoạn cố định và truy xuất dense thuần, chỉ \SL{retrieval_grid/grid/fixed/dense/hit@1Pct}\% câu hỏi có đoạn đúng ở vị trí
        đầu.
  \item \textbf{Câu hỏi đa bước không nằm trong đoạn nào.} ``Fed tăng lãi suất thì vàng SJC ra sao'' cần nối ít
        nhất ba mẩu tri thức nằm ở ba tài liệu khác nhau; top-$k$ theo độ tương đồng với câu hỏi không bảo đảm lấy
        đủ cả ba.
\end{enumerate}

\section{Thiết lập thực nghiệm truy xuất}\label{sec:thietlap}
\subsection{Kho tri thức và bộ câu hỏi}
Kho tri thức gồm \SL{retrieval_grid/n_docs} tài liệu Markdown, tám tài liệu y tế và chín tài liệu tài chính
(Bảng~\ref{tab:corpus}). Mỗi tài liệu tổng hợp từ một văn bản gốc có số hiệu, viết lại bằng lời của tác giả và có
cấu trúc tiêu đề hai cấp, mô phỏng cách một phác đồ hay một báo cáo phân tích được tổ chức trong thực tế.

\bang{tab_corpus}{Kho tri thức \SL{retrieval_grid/n_docs} tài liệu của tiểu luận}{corpus}

Bộ đánh giá có \SL{retrieval_grid/n_questions} câu hỏi (\SL{dem/qa_health} y tế, \SL{dem/qa_finance} tài chính).
Mỗi câu hỏi đi kèm một \emph{chuỗi đáp án} nằm nguyên văn trong kho (Bảng~\ref{tab:qa}); câu hỏi được viết bằng
cách diễn đạt khác với văn bản (``tiểu đường'' thay cho ``đái tháo đường'', ``uống bia'' thay cho ``rượu bia'') để
không chấm điểm sự trùng từ đơn thuần.

\bang{tab_qa_sample}{Trích bộ câu hỏi đánh giá truy xuất}{qa}

\subsection{Độ đo}
Một đoạn được coi là \textbf{liên quan} nếu nó chứa trọn chuỗi đáp án; định nghĩa này không cần gán nhãn tay và
công bằng giữa các chiến lược chia đoạn, vì đoạn bị cắt ngang đáp án không được tính. Với $r$ là thứ hạng của đoạn
liên quan đầu tiên trong danh sách trả về, ba độ đo được dùng (Hình~\ref{fig:dodo}):
\begin{align}
  \text{Hit@}k &= \mathbb{1}[r \le k], \qquad \text{MRR} = \frac{1}{|Q|}\sum_{q \in Q} \frac{1}{r_q},\\
  \text{nDCG@}10 &= \frac{1}{\text{IDCG@}10}\sum_{i=1}^{10}\frac{\text{rel}_i}{\log_2(i+1)},
\end{align}
trong đó $\text{rel}_i \in \{0,1\}$ cho biết đoạn ở vị trí $i$ có liên quan không và IDCG là DCG của thứ tự lý
tưởng. Hit@1 trả lời ``đoạn đầu tiên có đúng không'', câu hỏi quan trọng nhất khi chỉ đưa ít đoạn vào LLM; MRR
phạt nhẹ khi đoạn đúng ở vị trí 2--3; nDCG thưởng việc đặt \emph{mọi} đoạn liên quan lên trên.

\hinhtikz{t3_metrics}{Cách tính Hit@k, MRR và nDCG@10 trên một danh sách xếp hạng}{dodo}

\maNguon{retrieval__rank_metrics}{Tính Hit@k, MRR và nDCG@10 cho một câu hỏi}

\section{Các chiến lược phân mảnh (chunking)}\label{sec:chunking}
\subsection{Bài toán đánh đổi: độ dài đoạn và khoảng chồng lấn}
Chọn cỡ đoạn là chọn giữa hai loại lỗi. Đoạn \emph{ngắn} chính xác nhưng dễ cắt rời đáp án khỏi ngữ cảnh của nó;
đoạn \emph{dài} giữ trọn ngữ cảnh nhưng vector của nó là trung bình của nhiều ý, nên khó khớp với một câu hỏi cụ
thể, và mỗi đoạn đưa vào LLM tốn nhiều token hơn. Chồng lấn (overlap) $o$ giữa hai đoạn kề giảm xác suất cắt
ngang câu: với cửa sổ $w$, bước trượt là $w-o$, số đoạn của tài liệu dài $L$ là $\lceil (L-o)/(w-o) \rceil$.

Thí nghiệm quét bảy cỡ đoạn từ 120 tới 1300 ký tự (chia đệ quy, chồng lấn 15\%) cho thấy cả hai mặt của đánh đổi
(Hình~\ref{fig:ret_sweep}, Bảng~\ref{tab:sweep}). Ở cỡ 120 ký tự, chỉ
\SL{retrieval_sweep/0/intactPct}\% câu hỏi còn đoạn chứa trọn đáp án, và dense Hit@1 chỉ đạt
\SL{retrieval_sweep/0/dense_hit@1Pct}\%. Tăng cỡ đoạn thì Hit@1 tăng, nhưng ngữ cảnh top-5 phình từ
\SL{retrieval_sweep/0/ctx_tokens_top5} lên \SL{retrieval_sweep/6/ctx_tokens_top5} token. Lưu ý một hiệu ứng
của cách chấm: ở cỡ 1300 ký tự kho chỉ còn \SL{retrieval_sweep/6/n_chunks} đoạn, gần với truy xuất cả tài liệu,
nên tìm đúng ``đoạn'' dễ hơn nhiều; Hit@1 cao ở đó không có nghĩa LLM nhận được ngữ cảnh gọn. Vùng 400--900 ký
tự là điểm cân bằng, và là cỡ được chọn cho các thí nghiệm sau.

\hinhpdf{ret_sweep}{1.0}{Quét cỡ đoạn: tỉ lệ trúng tăng theo cỡ đoạn trong khi số ký tự ngữ cảnh top-5 tăng tuyến
tính}{ret_sweep}

\bang{tab_sweep}{Ảnh hưởng của cỡ đoạn tới truy xuất BM25 và dense}{sweep}

\subsection{Năm chiến lược phân mảnh}
Năm chiến lược được cài trong \texttt{tl2/chunking.py}, mọi chiến lược đều giữ vị trí ký tự gốc của đoạn để chấm
điểm được (Hình~\ref{fig:chunking}):

\begin{enumerate}
  \item \textbf{Cố định (fixed-size)}: cửa sổ trượt 400 ký tự, chồng lấn 60, không quan tâm ranh giới câu.
  \item \textbf{Đệ quy (recursive)}: thử cắt ở dấu phân cách cấp cao nhất (đoạn trống, xuống dòng, dấu chấm,
        chấm phẩy, dấu phẩy, khoảng trắng) và chỉ xuống cấp thấp hơn khi mảnh vẫn quá dài; các mảnh nhỏ liền nhau
        được gộp tới giới hạn.
  \item \textbf{Theo cấu trúc (structure-aware)}: mỗi mục Markdown là một đoạn, gắn chuỗi tiêu đề làm ngữ cảnh
        (``Thuốc điều trị > Metformin > Chống chỉ định''); mục quá 900 ký tự mới tách đệ quy bên trong.
  \item \textbf{Ngữ nghĩa (semantic)}: nhúng từng câu, cắt khi khoảng cách cosin giữa hai câu liền nhau vượt phân
        vị thứ 80 của chuỗi khoảng cách trong tài liệu.
  \item \textbf{Cha -- con (parent-child, small-to-big)}: lập chỉ mục các đoạn con khoảng 220 ký tự để khớp chính
        xác, nhưng trả cho LLM đoạn cha là cả mục (Hình~\ref{fig:chacon}).
\end{enumerate}

\hinhtikz{t3_chunking}{Năm chiến lược chia đoạn trên cùng một tài liệu: dải ngang là tài liệu, mỗi khối màu là
một đoạn}{chunking}

\hinhtikz{t3_parent_child}{Truy xuất cha -- con: khớp trên đoạn con nhỏ, đưa đoạn cha đủ ngữ cảnh cho LLM}{chacon}

\maNguon{chunking___recursive_spans}{Lõi của chia đoạn đệ quy: cắt ở dấu phân cách cấp cao nhất có thể}
\maNguon{chunking__semantic}{Chia đoạn ngữ nghĩa theo phân vị khoảng cách cosin giữa các câu liền nhau}
\maNguon{chunking__parent_child}{Chia đoạn cha -- con: cha là mục, con là mảnh đệ quy nhỏ bên trong}

Bảng~\ref{tab:chunkstats} và Hình~\ref{fig:ret_chunk_lengths} cho thấy hình dạng các đoạn sinh ra. Cố định cho độ
dài gần như đều; theo cấu trúc dao động theo độ dài mục; ngữ nghĩa dao động mạnh nhất, có đoạn chỉ
\SL{retrieval_grid/chunk_stats/semantic/min_len} ký tự (một dòng tiêu đề bị tách riêng) và có đoạn tới
\SL{retrieval_grid/chunk_stats/semantic/max_len} ký tự. Chia cố định là chiến lược duy nhất làm mất đáp án: chỉ
\SL{retrieval_grid/chunk_stats/fixed/intactPct}\% câu hỏi còn đoạn chứa trọn chuỗi đáp án.

\bang{tab_chunk_stats}{Thống kê đoạn của năm chiến lược trên toàn kho}{chunkstats}

\hinhpdf{ret_chunk_lengths}{1.0}{Phân bố độ dài đoạn của năm chiến lược}{ret_chunk_lengths}

Hình~\ref{fig:ret_semantic} mở hộp chiến lược ngữ nghĩa trên tài liệu \texttt{h03-ruou-bia}: mỗi cột là khoảng
cách giữa hai câu liền nhau, đường ngang là ngưỡng phân vị 80, cột đỏ là điểm cắt. Các điểm cắt rơi đúng vào chỗ
văn bản chuyển chủ đề (từ cơ chế nhiễm toan lactic sang nguy cơ hạ glucose huyết), nhưng cũng có điểm cắt giữa
một danh sách gạch đầu dòng vì các mục của danh sách ngắn và khác từ vựng.

\hinhpdf{ret_semantic}{1.0}{Chia đoạn ngữ nghĩa: khoảng cách cosin giữa các câu liền nhau và điểm cắt}{ret_semantic}

\subsection{Ý nghĩa của chunking trong dữ liệu y tế và tài chính Việt Nam}
Văn bản y tế và tài chính có hai đặc điểm làm chunking quan trọng hơn văn bản thông thường. Thứ nhất, \emph{con
số gắn chặt với ngữ cảnh}: ``$\ge 6{,}5\%$'' chỉ có nghĩa khi đi cùng ``HbA1c'' và ``tiêu chuẩn chẩn đoán''; một
đoạn chứa con số nhưng mất chủ ngữ là một đoạn nguy hiểm. Thứ hai, \emph{tài liệu có cấu trúc chuẩn}: phác đồ Bộ
Y tế, nghị định, báo cáo thường niên đều chia mục rõ ràng. Chiến lược theo cấu trúc khai thác đúng đặc điểm thứ
hai để bảo vệ đặc điểm thứ nhất: tiêu đề mục được gắn vào đoạn, nên đoạn ``Chống chỉ định'' luôn biết nó thuộc
mục ``Metformin''. Tiếng Việt còn có thêm một khó khăn: từ được ghép từ nhiều âm tiết cách nhau bởi khoảng trắng,
nên cắt theo khoảng trắng có thể tách ``đái tháo'' khỏi ``đường''; chia đệ quy ưu tiên dấu câu trước khoảng trắng
để tránh điều này.

\section{Truy xuất: thưa, dày, lai và xếp lại}
\subsection{Truy xuất thưa: TF-IDF và BM25}
Truy xuất thưa biểu diễn văn bản bằng tần suất từ. BM25~\cite{robertson2009bm25} cho điểm tài liệu $d$ với câu hỏi
$q$ gồm các từ $t$ bằng
\begin{equation}
  \operatorname{BM25}(q,d) = \sum_{t \in q} \operatorname{IDF}(t)\,
  \frac{f(t,d)\,(k_1+1)}{f(t,d) + k_1\left(1 - b + b\,\frac{|d|}{\operatorname{avgdl}}\right)},
\end{equation}
với $f(t,d)$ là tần suất từ, $|d|$ là độ dài đoạn, $k_1 = 1{,}5$ điều tiết bão hòa tần suất và $b = 0{,}75$ chuẩn
hóa độ dài. Tách từ ở mức âm tiết viết thường, giữ số và đơn vị như ``mg'', ``dl''.

\maNguon{retrieval__BM25Retriever}{Bộ truy xuất BM25 trên các âm tiết tiếng Việt}

\subsection{Truy xuất dày: mô hình nhúng đa ngôn ngữ}
Truy xuất dày~\cite{karpukhin2020dpr} dùng một bộ mã hóa Transformer ánh xạ câu hỏi và đoạn văn vào cùng một không
gian vector. Bài dùng \texttt{multilingual-e5-base}~\cite{wang2024me5}, mô hình được huấn luyện tương phản trên dữ
liệu đa ngôn ngữ, yêu cầu tiền tố \texttt{query:} cho câu hỏi và \texttt{passage:} cho đoạn văn. Đây là kiến trúc
\textbf{bi-encoder}~\cite{reimers2019sbert}: câu hỏi và đoạn văn được mã hóa \emph{riêng}, nên vector đoạn văn tính trước một lần khi lập
chỉ mục (Hình~\ref{fig:bicross}).

\maNguon{retrieval__Embedder}{Bọc SentenceTransformer cho mô hình e5: tiền tố, chuẩn hóa L2, chạy GPU}

So sánh ba mô hình nhúng (Bảng~\ref{tab:embed}, Hình~\ref{fig:ret_embedders}) cho thấy họ e5 vượt xa
\texttt{paraphrase-multilingual-MiniLM}: Hit@1 \SL{retrieval_embed/e5-base/hit@1Pct}\% so với
\SL{retrieval_embed/minilm/hit@1Pct}\%. Bản e5-small và e5-base có Hit@1 bằng nhau trên bộ câu hỏi này nhưng phân
bổ khác theo miền: e5-base tốt hơn ở y tế, e5-small tốt hơn ở tài chính. Bài chọn e5-base vì y tế là miền đòi hỏi
độ chính xác cao hơn.

\bang{tab_embed}{So sánh ba mô hình nhúng đa ngôn ngữ (truy xuất dense, đoạn theo cấu trúc)}{embed}

\hinhpdf{ret_embedders}{1.0}{Ba mô hình nhúng đa ngôn ngữ trên cùng bộ câu hỏi}{ret_embedders}

Hình~\ref{fig:ret_pca} chiếu \SL{retrieval_grid/chunk_stats/structure/n_chunks} đoạn theo cấu trúc xuống hai
chiều bằng PCA: đoạn y tế và đoạn tài chính tách thành hai vùng rõ rệt, và trong mỗi vùng các đoạn cùng tài liệu
nằm gần nhau. Hai thành phần đầu chỉ giữ \SL{retrieval_viz/embedding_pca/var_ratio/0Pct}\% và
\SL{retrieval_viz/embedding_pca/var_ratio/1Pct}\% phương sai, nhắc rằng không gian 768 chiều phức tạp hơn nhiều so
với những gì một hình hai chiều thể hiện. Hình~\ref{fig:ret_simmatrix} là ma trận cosin giữa mười câu hỏi và các
đoạn: ô sáng nhất mỗi hàng thường là đoạn đúng, nhưng khoảng cách giữa ô đúng và các ô kế bên nhỏ, lý do cần thêm
bước xếp lại.

\hinhpdf{ret_pca}{0.92}{Các đoạn trong không gian nhúng e5-base chiếu PCA hai chiều}{ret_pca}
\hinhpdf{ret_simmatrix}{1.0}{Độ tương đồng cosin giữa câu hỏi và đoạn văn}{ret_simmatrix}

\subsection{Truy xuất lai bằng Reciprocal Rank Fusion}
BM25 và dense mắc lỗi ở những câu hỏi khác nhau, nên gộp chúng lại có lợi. Điểm số của hai bộ truy xuất không
cùng thang đo (BM25 không giới hạn, cosin trong $[-1,1]$), nên thay vì cộng điểm, Reciprocal Rank Fusion
(RRF)~\cite{cormack2009rrf} cộng nghịch đảo thứ hạng:
\begin{equation}
  \operatorname{RRF}(d) = \sum_{r \in \{\text{BM25},\ \text{dense}\}} \frac{1}{k + \operatorname{rank}_r(d)}, \qquad k = 60.
\end{equation}
Hằng số $k$ làm phẳng đóng góp của các vị trí đầu, để một đoạn đứng hạng 1 ở một bộ truy xuất nhưng hạng 50 ở bộ
kia không áp đảo một đoạn đứng hạng 3 ở cả hai.

\maNguon{retrieval__HybridRRF}{Truy xuất lai: hợp nhất thứ hạng BM25 và dense bằng RRF}

\subsection{Xếp lại bằng cross-encoder}
Bi-encoder nhanh nhưng không bao giờ nhìn câu hỏi và đoạn văn \emph{cùng lúc}. Cross-encoder~\cite{nogueira2019rerank}
nối hai chuỗi thành một đầu vào và cho Transformer chú ý chéo giữa chúng, đầu ra là một điểm liên quan trực tiếp
$s(q,d) = \sigma\bigl(\mathbf{w}^\top \mathrm{Enc}([q; d])_{\texttt{[CLS]}}\bigr)$ (Hình~\ref{fig:bicross}). Cái
giá là phải chạy mô hình một lần cho \emph{mỗi} cặp, nên chỉ dùng được ở giai đoạn hai trên một tập ứng viên nhỏ.
Bài dùng \texttt{bge-reranker-v2-m3}~\cite{chen2024bgem3}, xếp lại 30 ứng viên của truy xuất lai
(Hình~\ref{fig:hybrid}).

\hinhtikz{t3_bi_cross}{Bi-encoder mã hóa riêng câu hỏi và đoạn văn rồi so cosin; cross-encoder mã hóa chung cả cặp
và cho điểm trực tiếp}{bicross}

\hinhtikz{t3_hybrid}{Phễu truy xuất lai: BM25 và dense chạy song song, hợp nhất RRF, cross-encoder xếp lại 30 ứng
viên và giữ top-$k$}{hybrid}

\maNguon{retrieval__Reranker}{Cross-encoder bge-reranker-v2-m3 chạy bf16 trên GPU}

\subsection{Kết quả: chiến lược chia đoạn $\times$ bộ truy xuất}
Bảng~\ref{tab:gridhit1} và Hình~\ref{fig:ret_grid_hit1} là kết quả chính của chương: 25 tổ hợp chia đoạn và truy
xuất trên \SL{retrieval_grid/n_questions} câu hỏi. Tổ hợp tốt nhất là \textbf{chia theo cấu trúc + truy xuất lai +
xếp lại}, Hit@1 \SL{retrieval_grid/grid/structure/hybrid_rerank/hit@1Pct}\% và MRR
\SL{retrieval_grid/grid/structure/hybrid_rerank/mrr}; tổ hợp Naive RAG (cố định + dense) chỉ đạt
\SL{retrieval_grid/grid/fixed/dense/hit@1Pct}\%.

\bang{tab_grid_hit1}{Hit@1 của 25 tổ hợp chiến lược chia đoạn và bộ truy xuất}{gridhit1}
\hinhpdf{ret_grid_hit1}{1.0}{Bản đồ nhiệt Hit@1: chiến lược chia đoạn (hàng) $\times$ bộ truy xuất (cột)}{ret_grid_hit1}
\bang{tab_grid_mrr}{MRR của 25 tổ hợp chiến lược chia đoạn và bộ truy xuất}{gridmrr}

Ba quan sát rút ra từ lưới:
\begin{itemize}
  \item \textbf{Chiến lược chia đoạn quan trọng ngang với bộ truy xuất.} Cùng bộ hybrid + rerank, đổi từ chia cố
        định sang chia theo cấu trúc nâng Hit@1 từ \SL{retrieval_grid/grid/fixed/hybrid_rerank/hit@1Pct}\% lên
        \SL{retrieval_grid/grid/structure/hybrid_rerank/hit@1Pct}\%.
  \item \textbf{Mỗi tầng của phễu đều có ích.} Với chia theo cấu trúc: BM25
        \SL{retrieval_grid/grid/structure/bm25/hit@1Pct}\%, dense \SL{retrieval_grid/grid/structure/dense/hit@1Pct}\%,
        lai \SL{retrieval_grid/grid/structure/hybrid/hit@1Pct}\%, lai + xếp lại
        \SL{retrieval_grid/grid/structure/hybrid_rerank/hit@1Pct}\%.
  \item \textbf{Cha -- con mạnh nhất khi không có reranker.} Khớp trên đoạn con nhỏ cho dense Hit@1
        \SL{retrieval_grid/grid/parent_child/dense/hit@1Pct}\%, cao nhất cột dense; khi đã có reranker thì ưu thế
        này thu hẹp.
\end{itemize}

Bảng~\ref{tab:retdetail} đi sâu vào chiến lược theo cấu trúc. Hình~\ref{fig:ret_hitk} là đường Hit@$k$: mọi bộ
truy xuất đều tìm được đoạn đúng trong top-10 ở gần như mọi câu hỏi; khác biệt nằm ở \emph{thứ hạng}, đúng điều
mà reranker cải thiện. Hình~\ref{fig:ret_domain} tách theo miền: câu hỏi y tế khó hơn tài chính với mọi bộ truy
xuất, vì nhiều tài liệu y tế cùng nói về một bệnh và dùng chung từ vựng (metformin xuất hiện trong bốn tài liệu).

\bang{tab_ret_detail}{Chi tiết năm bộ truy xuất trên đoạn chia theo cấu trúc}{retdetail}
\hinhpdf{ret_hitk}{0.95}{Đường Hit@$k$ của năm bộ truy xuất}{ret_hitk}
\hinhpdf{ret_domain}{0.95}{Hit@1 theo miền y tế và tài chính}{ret_domain}

Chất lượng có giá. Hình~\ref{fig:ret_latency} đặt Hit@1 cạnh độ trễ: BM25 trả lời trong
\SL{retrieval_grid/grid/structure/bm25/latency_ms} ms, dense trong \SL{retrieval_grid/grid/structure/dense/latency_ms}
ms (gồm mã hóa câu hỏi trên GPU), còn hybrid + rerank mất \SL{retrieval_grid/grid/structure/hybrid_rerank/latency_ms}
ms vì phải chạy cross-encoder 30 lần. Với một tác tử hội thoại, nửa giây là chấp nhận được; với tìm kiếm thời gian
thực trên hàng triệu đoạn thì cần giảm số ứng viên hoặc dùng reranker nhỏ hơn.

\hinhpdf{ret_latency}{0.95}{Chất lượng và độ trễ của 25 tổ hợp: xếp lại đắt hơn nhiều nhưng chính xác hơn}{ret_latency}

\section{Advanced RAG: tối ưu đa chặng}
Advanced RAG giữ khung truy xuất -- sinh của Naive RAG nhưng thêm bước tối ưu ở hai đầu:

\begin{itemize}
  \item \textbf{Tiền truy xuất} sửa \emph{câu hỏi}: viết lại câu hỏi cho rõ, tách câu hỏi phức thành câu hỏi con,
        sinh nhiều cách hỏi (Multi-Query), hoặc sinh một câu trả lời giả định rồi dùng nó để tìm (HyDE).
  \item \textbf{Truy xuất} cải tiến bằng chia đoạn tốt hơn, truy xuất lai, lọc theo metadata (miền, tài liệu).
  \item \textbf{Hậu truy xuất} sửa \emph{ngữ cảnh}: xếp lại bằng cross-encoder, nén hoặc cắt bớt đoạn, sắp xếp lại
        để thông tin quan trọng nằm ở đầu và cuối prompt.
\end{itemize}

Phần truy xuất lai và xếp lại đã được đo ở mục trước. Mục này đo hai kỹ thuật tiền truy xuất dùng LLM cục bộ
(Hình~\ref{fig:hyde}).

\textbf{HyDE} (Hypothetical Document Embeddings)~\cite{gao2023hyde} dựa trên nhận xét: câu hỏi và đoạn trả lời
thường không giống nhau về hình thức, nhưng hai đoạn trả lời thì giống nhau. Thay vì nhúng câu hỏi $q$, HyDE cho
LLM viết một đoạn trả lời giả định $\tilde{d} = \mathrm{LLM}(q)$ rồi tìm bằng $f(\tilde{d})$. Biến thể
\emph{HyDE + câu hỏi} hợp nhất bằng RRF hai danh sách tìm bằng $f(q)$ và $f(\tilde d)$.
\textbf{Multi-Query} cho LLM viết ba cách hỏi khác $q_1, q_2, q_3$, truy xuất cho từng câu và hợp nhất bốn danh
sách bằng RRF.

\hinhtikz{t3_hyde}{Hai kỹ thuật tối ưu tiền truy xuất bằng LLM cục bộ: HyDE tìm bằng câu trả lời giả định,
Multi-Query hợp nhất nhiều cách hỏi bằng RRF}{hyde}

\maNguon{run_advanced__rrf}{Hợp nhất nhiều danh sách xếp hạng bằng RRF trong thí nghiệm Advanced RAG}

Kết quả (Bảng~\ref{tab:advanced}, Hình~\ref{fig:adv_rag}) là một kết quả âm có giá trị. Với Qwen2.5-1.5B, HyDE
\emph{làm giảm} Hit@1 từ \SL{advanced/agg/dense/hit@1Pct}\% xuống \SL{advanced/agg/hyde/hit@1Pct}\%, thua dense gốc
ở \SL{advanced/agg/hyde/losses_vs_dense} câu và chỉ thắng ở \SL{advanced/agg/hyde/wins_vs_dense} câu.
Multi-Query giảm ít hơn (\SL{advanced/agg/multiquery/hit@1Pct}\%) nhưng vẫn giảm. Bảng~\ref{tab:hydeex} giải thích
lý do: đoạn giả định của mô hình nhỏ chứa các chi tiết \emph{bịa}, ví dụ nói metformin cùng rượu gây ``tăng huyết
áp'', và vector của chi tiết bịa kéo truy xuất về các đoạn sai. HyDE chỉ có lợi khi LLM sinh đoạn giả định đủ đúng
để trông giống đoạn thật; với một mô hình 1,5 tỉ tham số trên miền chuyên sâu tiếng Việt, điều kiện đó không thỏa.
Mỗi câu hỏi còn tốn thêm \SL{advanced/gen_s/hyde} s sinh văn bản. Vì vậy tác tử ở Chương~\ref{chap:agent} không
dùng HyDE hay Multi-Query mà dùng truy xuất lai + xếp lại.

\bang{tab_advanced}{Tối ưu tiền truy xuất bằng LLM so với truy xuất dense gốc}{advanced}
\hinhpdf{adv_rag}{1.0}{Ảnh hưởng của HyDE và Multi-Query: Hit@1, MRR và số câu thắng, thua so với dense gốc}{adv_rag}
\bang{tab_hyde_ex}{Ví dụ văn bản do Qwen2.5-1.5B sinh trong HyDE và Multi-Query}{hydeex}

\section{GraphRAG: tổng hợp tri thức trên đồ thị}
\subsection{Ý tưởng và quy trình}
Advanced RAG vẫn coi tri thức là một tập đoạn văn rời rạc. GraphRAG~\cite{edge2024graphrag} thay đơn vị tri thức
bằng \textbf{đồ thị thực thể}, rồi thêm một tầng trừu tượng là \textbf{cộng đồng}: nhóm thực thể liên kết chặt
được tóm tắt trước thành một ``báo cáo cộng đồng''. Quy trình gồm bốn bước (Hình~\ref{fig:graphrag}):

\begin{enumerate}
  \item \textbf{Đồ thị thực thể}: ở bài báo gốc, LLM trích xuất thực thể và quan hệ từ văn bản; ở đây hai KG đã có
        sẵn và được hợp nhất thành một đồ thị vô hướng có trọng số (cạnh nhiều quan hệ được cộng trọng số).
  \item \textbf{Phát hiện cộng đồng}: thuật toán Louvain~\cite{blondel2008louvain} tìm phân hoạch $\{c_i\}$ cực đại
        độ đo modularity
        \begin{equation}
          Q = \frac{1}{2m}\sum_{i,j}\Bigl(A_{ij} - \frac{k_i k_j}{2m}\Bigr)\delta(c_i, c_j),
        \end{equation}
        với $A_{ij}$ là trọng số cạnh, $k_i$ là bậc có trọng số, $m$ là tổng trọng số. $Q$ đo mật độ cạnh trong cộng
        đồng vượt bao nhiêu so với một đồ thị ngẫu nhiên cùng dãy bậc.
  \item \textbf{Báo cáo cộng đồng}: các bộ ba nằm trọn trong một cộng đồng được diễn đạt thành câu và LLM viết một
        đoạn tóm tắt 3--4 câu.
  \item \textbf{Tìm kiếm}: \emph{cục bộ} liên kết thực thể trong câu hỏi rồi lấy các bộ ba quanh chúng và đoạn
        văn của tài liệu nguồn; \emph{toàn cục} xếp hạng các báo cáo cộng đồng theo độ tương đồng với câu hỏi.
\end{enumerate}

\hinhtikz{t3_graphrag}{Đường ống GraphRAG: đồ thị thực thể, cộng đồng Louvain, báo cáo cộng đồng do LLM viết,
tìm kiếm cục bộ và toàn cục}{graphrag}

Bước khó nhất trong thực tế là \textbf{liên kết thực thể}: nhận ra ``đang uống glucophage'' nói về
\textsf{Metformin}. Bộ liên kết của bài (Đoạn mã~\ref{code:graphrag__EntityLinker}) bỏ dấu tiếng Việt, khớp nguyên
từ với bảng bí danh và ưu tiên bí danh dài nhất trước, để ``gà rán'' không bị hiểu thành ``gà'' và một đoạn đã
khớp không bị khớp lại.

\maNguon{graphrag__EntityLinker}{Liên kết thực thể: khớp bí danh dài nhất trước, không dấu, nguyên từ}
\maNguon{graphrag__summarize_communities}{Tóm tắt cộng đồng bằng LLM cục bộ từ các bộ ba nằm trong cộng đồng}
\maDai{graphrag__GraphRAG}{Lớp \texttt{GraphRAG}: tìm kiếm cục bộ trên KG và tìm kiếm toàn cục trên báo cáo cộng đồng}

\subsection{Cộng đồng tri thức}
Louvain tìm được \SL{graphrag/n_comms} cộng đồng với modularity \SL{graphrag/modularity}
(Hình~\ref{fig:grag_communities}, Bảng~\ref{tab:communities}). Không cộng đồng nào chứa thực thể của cả hai miền,
đúng như mong đợi vì hai KG không có cạnh nối nhau. Các cộng đồng có nghĩa rõ ràng mà không cần ai gán nhãn: C0
gom đái tháo đường típ 2 với tiêu chuẩn chẩn đoán, yếu tố nguy cơ, các món tinh bột và lối sống; C2 gom ba thuốc
chính (metformin, insulin, gliclazide) với rượu bia, nhiễm toan lactic, hạ glucose huyết và các chống chỉ định,
tức đúng phân hệ ``an toàn thuốc''; C6 là bộ ba empagliflozin -- suy tim -- bệnh thận mạn, đúng chỉ định ưu tiên
của nhóm SGLT2. Bên tài chính, C1 là kênh tỷ giá -- lãi suất -- cổ phiếu và ngành, C3 là kênh vàng thế giới (Fed,
DXY, XAU/USD), C4 là chính sách vàng miếng trong nước. LLM viết \SL{dem/comms} bản tóm tắt trong
\SL{graphrag/summary_s} s.

\hinhpdf{grag_communities}{1.0}{Các cộng đồng Louvain trên đồ thị hợp nhất hai KG; nhãn ghi kích thước và thực thể
trung tâm, cạnh đứt nối hai cộng đồng khác nhau}{grag_communities}

\bang{tab_communities}{Các cộng đồng tri thức và báo cáo cộng đồng}{communities}

\subsection{So sánh bốn kiến trúc trên câu hỏi đa bước}
Để đo đúng thế mạnh của GraphRAG, bài dùng một bộ \SL{dem/grag_q} câu hỏi đa bước riêng, mỗi câu kèm các nhóm từ
khóa ứng với các dữ kiện bắt buộc phải có (ví dụ câu Fed -- vàng SJC cần cả kênh DXY, kênh XAU, kênh tỷ giá và
phần bù trong nước). Bốn kiến trúc dùng cùng LLM và cùng prompt có trích dẫn, chỉ khác ngữ cảnh:
Naive (cố định + dense, top-5), Advanced (cấu trúc + lai + xếp lại, top-5), GraphRAG cục bộ (bộ ba quanh thực thể
liên kết + 3 đoạn văn của tài liệu nguồn) và GraphRAG cục bộ + toàn cục (thêm 2 báo cáo cộng đồng gần nhất).
\textbf{Độ phủ ngữ cảnh} là tỉ lệ dữ kiện có mặt trong ngữ cảnh đưa vào LLM; \textbf{độ phủ câu trả lời} là tỉ lệ dữ
kiện có trong câu trả lời cuối.

\bang{tab_graphrag}{Bốn kiến trúc RAG trên \SL{dem/grag_q} câu hỏi đa bước}{graphrag}
\hinhpdf{grag_coverage}{1.0}{Độ phủ dữ kiện trong ngữ cảnh và trong câu trả lời, cùng độ dài ngữ cảnh của bốn kiến
trúc}{grag_coverage}

GraphRAG cục bộ cho câu trả lời đầy đủ nhất: độ phủ câu trả lời \SL{graphrag/agg/graph_local/ans_covPct}\% so với
\SL{graphrag/agg/naive/ans_covPct}\% của Naive và \SL{graphrag/agg/advanced/ans_covPct}\% của Advanced. Điều đáng
chú ý là Advanced có độ phủ \emph{ngữ cảnh} cao nhất (\SL{graphrag/agg/advanced/ctx_covPct}\%) nhưng độ phủ câu
trả lời thấp hơn GraphRAG: dữ kiện có trong năm đoạn văn dài, nhưng LLM nhỏ không gom đủ chúng lại. Bộ ba KG đã
diễn đạt thành câu ngắn, mỗi câu một quan hệ, nên LLM dễ chép đúng vào câu trả lời hơn. Thêm báo cáo toàn cục đưa
độ phủ ngữ cảnh lên \SL{graphrag/agg/graph_local_global/ctx_covPct}\% nhưng \emph{giảm} độ phủ câu trả lời xuống
\SL{graphrag/agg/graph_local_global/ans_covPct}\%, vì ngữ cảnh dài thêm gần gấp rưỡi
(\SL{graphrag/agg/graph_local_global/ctx_chars} ký tự) và mô hình nhỏ bị loãng sự chú ý. Tìm kiếm toàn cục hợp
với câu hỏi tổng hợp kiểu ``chủ đề chính của kho là gì'', không phải câu hỏi đa bước có đáp án cụ thể.

\bang{tab_grag_perq}{Độ phủ ngữ cảnh / câu trả lời theo từng câu hỏi đa bước}{gragperq}

\section{Chiến lược lựa chọn kiến trúc trong hệ thống thực tế}
Không có kiến trúc RAG tốt nhất cho mọi trường hợp. Bảng~\ref{tab:chonkientruc} tóm tắt hướng dẫn chọn rút ra từ
các thí nghiệm của chương.

\begin{table}[H]\centering
\caption{Hướng dẫn chọn kiến trúc RAG theo loại câu hỏi và điều kiện triển khai}\label{tab:chonkientruc}
{\small\begin{tabular}{L{4.4cm}L{4.4cm}L{5.4cm}}
\toprule
Tình huống & Kiến trúc nên dùng & Căn cứ trong bài \\
\midrule
Hỏi đáp tra cứu một dữ kiện & Advanced: cấu trúc + lai + xếp lại & Hit@1 cao nhất trên \SL{retrieval_grid/n_questions} câu \\
Không có GPU, cần độ trễ thấp & BM25 hoặc lai không xếp lại & nhanh hơn reranker hàng trăm lần \\
Câu hỏi đa bước, nhân quả & GraphRAG cục bộ & độ phủ câu trả lời cao nhất \\
Câu hỏi tổng hợp toàn kho & GraphRAG toàn cục & báo cáo cộng đồng tóm tắt sẵn \\
Câu hỏi có điều kiện số & truy vấn KG / SPARQL & TF-IDF bỏ sót gần hết (Mục 2.4) \\
LLM nhỏ ($\le$ 2 tỉ tham số) & tránh HyDE, Multi-Query & cả hai làm giảm Hit@1 \\
\bottomrule
\end{tabular}}
\end{table}

Tác tử của bài kết hợp ba hàng của bảng: truy xuất lai + xếp lại cho đoạn văn, KG cho quan hệ và ràng buộc, tính
toán trực tiếp cho số liệu thị trường, và để bộ lập kế hoạch chọn công cụ theo loại câu hỏi (Chương~\ref{chap:agent}).
